\clearpage
\section{기약 사차식의 다섯 갈루아군}
\label{sec:quartic-group-classification}

\begin{learninggoal}
$S_4$의 추이적 부분군을 분류하고, 삼차분해식의 인수분해형과 판별식으로 $V_4,C_4,D_4,A_4,S_4$를 판정한다. $C_4$와 $D_4$를 분해체 차수와 제곱류로 구분한다.
\end{learninggoal}

$f\in K[X]$가 기약인 분리 사차다항식이면 갈루아군 $G=\Gal(f/K)$는 네 근 위에서 추이적으로 작용한다. 궤도-안정자 정리에 의해 $4$가 $|G|$를 나눈다.

\begin{theorem}[$S_4$의 추이적 부분군]
\label{thm:transitive-subgroups-s4}
켤레를 제외하면 $S_4$의 추이적 부분군은 정확히 다음 다섯 개이다.
\[
  V_4,\qquad C_4,\qquad D_4,\qquad A_4,\qquad S_4.
\]
그 위수는 각각
\[
  4,\qquad4,\qquad8,\qquad12,\qquad24
\]
이다.
\end{theorem}

\begin{proof}
앞 절의 전사준동형
\[
  \rho:S_4\to S_3,
  \qquad \ker\rho=V_4
\]
를 $G$에 제한하고 $N=G\cap V_4$라 두자. 그러면
\[
  G/N\cong\rho(G).
\]

먼저 $\rho(G)=1$이면 $G\subseteq V_4$이다. $G$가 추이적이려면 위수가 $4$의 배수여야 하므로 $G=V_4$이다.

$\rho(G)\cong C_3$이면 $G$는 $\rho^{-1}(A_3)=A_4$ 안에 들어간다. 한편 $4\mid |G|$이고 $3\mid |G|$이므로 $12\mid |G|$이다. 따라서 $G=A_4$이다.

$\rho(G)=S_3$이면 $6\mid |G|$이고 $4\mid |G|$이므로 $12\mid |G|$이다. 위수 $12$인 $S_4$의 부분군은 지수 $2$이므로 정규이고, 유일한 지수 $2$ 부분군은 $A_4$이다. 그러나 $A_4$의 상은 $A_3$이므로 모순이다. 따라서 $|G|=24$이고 $G=S_4$이다.

마지막으로 $\rho(G)\cong C_2$라 하자. 그러면
\[
  |G|=2|N|.
\]
추이성 때문에 $4\mid|G|$이고 $N\subseteq V_4$이므로 $|N|=2$ 또는 $4$이다. 첫 경우 $|G|=4$이고, $G$는 $V_4$가 아니므로 순환군 $C_4$이다. 둘째 경우 $|G|=8$이며 $S_4$의 추이적 위수 $8$ 부분군은 정사각형의 대칭군 $D_4$와 켤레이다.
\end{proof}

\subsection{삼차분해식과 판별식의 결정표}

$f=X^4+pX^2+qX+r$가 기약이고 분리이며 $R_f$가 앞 절의 삼차분해식이라 하자.

\begin{theorem}[기약 사차식의 판정]
\label{thm:quartic-galois-decision}
다음이 성립한다.
\begin{enumerate}[label=(\roman*)]
  \item $R_f$가 $K$에서 완전히 분해되면
  \[
    G\cong V_4.
  \]
  \item $R_f$가 기약이면
  \[
    G\cong
    \begin{cases}
      A_4,&\Disc(f)\in K^{\times2},\\
      S_4,&\Disc(f)\notin K^{\times2}.
    \end{cases}
  \]
  \item $R_f$가 일차인수와 기약 이차인수의 곱이면
  \[
    G\cong C_4\quad\text{또는}\quad D_4.
  \]
  이 경우 $F$를 $R_f$의 분해체라 하면
  \[
    G\cong C_4\iff[L:F]=2,
    \qquad
    G\cong D_4\iff[L:F]=4.
  \]
\end{enumerate}
\end{theorem}

\begin{proof}
정리 \ref{thm:quartic-resolvent-fixed-field}에서
\[
  \Gal(F/K)\cong\rho(G)
\]
이다. $R_f$의 인수분해형은 세 근 위의 작용의 궤도분해와 같으므로, 완전분해·기약·일차 곱하기 기약이차의 세 경우에 $\rho(G)$는 각각
\[
  1,\qquad C_3\text{ 또는 }S_3,\qquad C_2
\]
이다. 정리 \ref{thm:transitive-subgroups-s4}를 적용하면 가능한 군들이 나온다.

$R_f$가 기약일 때 $A_4$와 $S_4$는 판별식의 교대군 판정으로 구별된다. 마지막 경우에는
\[
  [L:F]=|G\cap V_4|
\]
이고 이 값이 $2$이면 $|G|=4$, $4$이면 $|G|=8$이다.
\end{proof}

\begin{center}
\renewcommand{\arraystretch}{1.25}
\begin{tabular}{lll}
\toprule
$R_f$의 인수분해형 & 추가 정보 & $G$ \\
\midrule
$1+1+1$ & 없음 & $V_4$ \\
$3$ & $\Disc(f)$ 제곱 & $A_4$ \\
$3$ & $\Disc(f)$ 비제곱 & $S_4$ \\
$1+2$ & $[L:F]=2$ & $C_4$ \\
$1+2$ & $[L:F]=4$ & $D_4$ \\
\bottomrule
\end{tabular}
\end{center}

\subsection{$C_4$와 $D_4$의 제곱류 판정}

$R_f=(Z-z_0)h(Z)$이고 $h$가 기약 이차라고 하자. $F$는 $h$의 분해체이므로 $[F:K]=2$이고, $R_f$의 세 근을 $z_0,z_2,z_3\in F$라 쓸 수 있다. 앞 절의 사차 공식에서
\[
  L=F\bigl(\sqrt{-z_0},\sqrt{-z_2},\sqrt{-z_3}\bigr).
\]
또한 적절한 부호를 택하면
\[
  \sqrt{-z_0}\sqrt{-z_2}\sqrt{-z_3}=-q.
\]

\begin{proposition}[실용적인 제곱 판정]
\label{prop:c4-d4-square-test}
위 상황에서 $q\ne0$이라 하자. 그러면
\[
  G\cong C_4
  \quad\Longleftrightarrow\quad
  -z_0\in F^{\times2},
\]
그리고 그렇지 않으면 $G\cong D_4$이다.
\end{proposition}

\begin{proof}
$F^\times/F^{\times2}$에서
\[
  [-z_0][-z_2][-z_3]=[q^2]=1.
\]
$-z_0$가 $F$의 제곱이면 $[-z_2]=[-z_3]$이다. 세 제곱근을 모두 넣는 데 하나의 독립적인 제곱류만 필요하므로 $[L:F]\le2$. $f$가 기약 사차식인데 $[F:K]=2$이므로 $L\ne F$이고, 따라서 $[L:F]=2$이다. 정리 \ref{thm:quartic-galois-decision}에서 $G\cong C_4$이다.

반대로 $G\cong C_4$이면 $[L:F]=2$이므로 세 제곱류가 생성하는 부분공간의 차원은 $1$이다. $F/K$의 비자명한 자동동형은 $z_2,z_3$을 맞바꾸고 $z_0$을 고정한다. 제곱인 원소는 제곱인 원소로 보내지므로 $[-z_2]$와 $[-z_3]$ 중 하나만 자명할 수는 없다. 곱이 자명한 세 제곱류가 한 차원 안에 놓이려면 $[-z_0]=1$이어야 한다.
\end{proof}

$q=0$이면 삼차분해식의 상수항이 $0$이고 $z_i$ 중 하나가 $0$이다. 이 경우에는 위 명제 대신
\[
  [L:F]=2\text{ 또는 }4
\]
를 직접 계산하는 편이 안전하다.

\begin{pitfall}
삼차분해식이 기약이라고 해서 사차식의 갈루아군이 자동으로 $S_4$인 것은 아니다. 판별식이 제곱이면 군은 $A_4$이다. 또한 삼차분해식에 한 근이 있다고 해서 군이 자동으로 $D_4$인 것도 아니다. 분해체 차수가 $4$인 순환 사차확장에서는 $C_4$가 나타난다.
\end{pitfall}

\begin{checkpoint}
\begin{enumerate}[label=(\alph*)]
  \item 기약 사차식의 삼차분해식이 완전히 분해되면 판별식이 자동으로 제곱인 이유를 설명하라.
  \item 삼차분해식이 기약이고 판별식이 제곱이 아니면 왜 $D_4$가 불가능한가?
  \item $R_f$가 $1+2$형일 때 $[L:K]$의 가능한 값을 구하라.
\end{enumerate}
\end{checkpoint}
